\subsection{The integral count at large dyadic conductor}
\label{sec:rh-dyadic-host-count}

The same graded length theorem applies to the original inertial
deformation ring; it is not necessary to replace it by a component
of a different ordinary deformation ring. We give the presentation,
the point conditions and the comparison at every coefficient modulus.

\begin{theorem}[The actual inertial length comparison]
\label{thm:rh-dyadic-host-count}
Keep the frame, coefficient ring $O$, mark $g_1$, and actual
universal representation $u$ of
Theorems~\ref{thm:rh-dyadic-host-frame}
and~\ref{thm:rh-dyadic-host-universal}.
No residual weight-minus-two vanishing is assumed.
There are finitely many original characteristic-zero type-I or
type-II representations $r_i$, pairwise nonisomorphic on $\ker\eta$,
and a nonnegative integer $V$ such that for every $0\ne a\in O$,
\[
 \operatorname{length}_O\operatorname{Sha}_O(a)
 \le V\le 2e_O\sum_i
 v_\ell\bigl(\operatorname{tr}r_i(\phi)
                 -\operatorname{tr}r_0(\phi)\bigr).
 \tag{DC1}
\]
Here $\phi$ is the single Frobenius element of
Theorem~\ref{thm:rh-dyadic-host-traces}, $e_O=v_\pi(\ell)$,
$\pi$ is a uniformizer of $O$, and $v_\ell(\ell)=1$.
Their restrictions to $G_K$ belong to the two classes of
Theorem~\ref{thm:support-spectrum}, for $K=\mathbf Q(i)$.
In particular the normalized strict Selmer length in the source
satisfies
\[
       (\log\ell)V_{\rm n}\le(\ell Ne)^C                         \tag{DC2}
\]
for an absolute exponent $C$ on these RH data.
\end{theorem}

\begin{proof}
Write $k=O/\pi O$, $F=\operatorname{Frac}O$, and use indices
$0,1,2$ for matrices. Set $T=\operatorname{diag}(1,0,-1)$,
let $J_3$ be the reversing permutation matrix, and put
$\theta(h)=J_3h^{-\mathsf t}J_3$.
The polarized group is $\mathrm{GL}_3\rtimes\{1,-1\}$,
with action $\theta$. The weights are the differences of the
diagonal entries of $T$. A normalized representation has matrix
one at $c$, and its two adjacent off-diagonal entries in
$r(g_1)-r(cg_1^{-1}c)$ are zero. In particular, writing only
matrix parts and $g^*=cg^{-1}c$, for $\eta(g)=1$ we have
\[
 r(g^*)=J_3r(g)^{\mathsf t}J_3,\qquad
 r(g_1)-r(g_1^*)=\zeta T,\qquad
 \bar\zeta=\overline{\Psi(g_1)-\Psi(g_1)^{-1}}\ne0.             \tag{DC3}
\]
The middle diagonal and the two corner entries of the difference
vanish by its antisymmetry under reversing transpose; the other
off-diagonal entries vanish by the two normalization equations.
These identities hold over local coefficient rings because $2$
is a unit.

\emph{Existence and minimal presentation.}
Here one can construct the original universal ring directly.
The kernel of the residual polarized representation is an open
normal subgroup of $\Gamma$. Its maximal pro-$\ell$ quotient is
topologically finitely generated: the first cohomology with
$\mathbf F_\ell$ coefficients of the corresponding number field
unramified outside the finite set of places is finite. Indeed the
Frattini quotient is dual to this cohomology, and lifts of a basis
generate the pro-$\ell$ group (a proper closed subgroup is
contained in a proper open subgroup, hence in a maximal one).
Quotienting $\Gamma$ by the kernel of this pro-$\ell$ quotient
therefore gives a topologically finitely generated profinite group.
Every residual deformation factors through it. For finite
coefficient quotients this follows from the congruence filtration,
whose successive kernels are additive groups of characteristic
$\ell$; for complete coefficients use reduction and separatedness.
Even with a general adically complete target, continuity gives
finite image at each discrete quotient.

Use matrices on finitely many topological generators with the
prescribed residual matrices as formal variables. Impose the group
relations, the multiplier, the two marks and the literal conditions
on all $N_v$ away from $\ell$. At $\ell$ the local element $x$
in (DH3) has distinct residual eigenvalues. Its matrix has a
canonical basis of ordered eigenvectors reducing to the standard
basis, obtained from the Hensel roots of its characteristic
polynomial and the spectral projectors, with each matching
coordinate normalized to one. An ordered flag congruent to the
standard one and stable under the multiplier-one subgroup is
necessarily the flag spanned by the successive initial segments
of that basis. Thus the existence of the original local conjugator
is exactly a set of equations: the indicated lower entries in
this basis vanish and the diagonal inertia entries are the
prescribed characters. In particular these are closed equations,
not a choice of additional local parameters. The resulting quotient
of the finite-variable power series ring represents the original
normalized functor, including on complete local rings with the
same residue field. All evaluation maps are local and continuous.
This construction gives its original $R,u$ by uniqueness.

Choose a minimal set of actual entry displacements forming a basis
of $\mathfrak m_R/(\mathfrak m_R^2+\pi R)$, as in (DU1);
write $n_j$ for the number of its variables of weight $j$.
They can be chosen homogeneously, since the torus acts on each
entry with its matrix weight. The formal torus argument in
Theorem~\ref{thm:rh-dyadic-host-universal} gives a weight-stable
kernel $\mathfrak a$ in
\[
                   P=O[[z_1,\ldots,z_m]]\twoheadrightarrow R .
\]
Every relation has constant coefficient zero by the distinguished
base point, and its linear coefficients belong to $\pi O$
by minimality. The only weights of variables are $-2,-1,0,1,2$.
Write $M_j$ for the corresponding residual adjoint module,
$H^1_L(j)$ for the original Selmer quotient of cocycles, and
$\operatorname{Ob}(j)$ for the original relative obstruction
space (the pair of a global two-cocycle and local primitives,
modulo global changes of lift and allowed local cocycles).
Put $h_j=\dim_k H^1_L(j)$ and $b_j=\dim_k\operatorname{Ob}(j)$.
Theorem~\ref{thm:rh-dyadic-host-frame}, in the normalization
(DH2)--(DH5), gives
\[
 h_j-b_j=
 \begin{cases}-1&j=-1,\\1&j=1,2,\\0&\text{otherwise}.\end{cases} \tag{DC4}
\]

We verify the local lifting and relation bound for this very
inertial functor. Consider a small extension of finite graded local
$O$-algebras $A'\to A$ with kernel $D$ killed by the maximal ideal
and concentrated in weight $j$. Matrices and all changes of lift
are taken homogeneously, so their multiplication defect is a
two-cocycle in $M_j\otimes_kD$.
Away from $\ell$, the original quotient is $D_v/N_v$:
its inertia has order prime to $\ell$ and its unramified quotient
is procyclic. Averaging on inertia and the cohomological dimension
of the procyclic quotient show that its positive-degree cohomology
in degrees at least two vanishes on these finite modules.
The local representation therefore lifts, still killing $N_v$.

At $\ell$, use the ordered eigenbasis of the local mark just
described. It is homogeneous: the characteristic-polynomial
coefficients and the uniquely specified Hensel roots have weight
zero; each spectral projector has the matrix weights. Applying
the projector to its corresponding standard vector and dividing
by the matching unit coordinate gives the required column.
The basis change lifts homogeneously to $A'$.
In this basis the matrix representation is upper triangular on
the multiplier-one subgroup. It is also upper triangular on the
other coset if that coset occurs. In fact $\theta$ preserves the
upper triangular group; for any element in that coset its
matrix intertwines the ordered eigenflags of $x$ and
$\theta(r(g^{-1}xg))$. Both ordered flags are standard, with
the matching residual ordering, so the intertwiner preserves
the flag. This argument holds over the finite local ring because
the residual eigenvalues are distinct.

First lift the diagonal representation, with exactly its given
inertia values. Dividing its diagonal characters by the base
diagonal representation leaves a cocycle of the unramified
procyclic quotient into diagonal units, with action through
$\eta$. Lifting the value on a Frobenius generator extends
continuously, since the coefficient groups at this step are
finite. This preserves all inertia values; diagonal entries
have weight zero. Lift the upper triangular entries arbitrarily
with this diagonal. Their defect now lies in the strictly upper
adjoint module of the residual representation, in weight $j$.
For $j>0$ its second cohomology is zero by \texttt{LocVan}
proved in Theorem~\ref{thm:rh-dyadic-host-frame}; for $j\le0$
the defect is zero already. Correct by an upper triangular
cochain, giving the original local lift. Conjugating back by the
lifted basis preserves the weights. Thus at every place there
are local lifts, with the same original local conditions.
There are no other exceptional places here: the finite character
has order prime to $\ell$.

Choose a continuous set-theoretic global lift and local lifts.
Its multiplication defect and their differences form exactly
the pair defining $\operatorname{Ob}(j)\otimes_k D$.
Changing the global lift adds a global coboundary and its
restrictions. The allowable local changes are precisely the
original residual local cocycles: at $v\ne\ell$ they vanish
on $N_v$; at $\ell$ a negative-weight change is a coboundary,
a zero-weight change vanishes on inertia, and a positive-weight
change is arbitrary. They preserve the local condition:
in the small kernel the integral basis change reducing to one
acts trivially, a negative coboundary is effected by an
infinitesimal conjugation, and the other two changes preserve
the upper triangular flag and its specified inertia diagonal.
Consequently vanishing of this relative obstruction suffices for
a global lift satisfying the original local equations.

It also suffices for a normalized lift. At $c$, the error in the
matrix identity is a matrix $X$ in the kernel satisfying
$X+d\theta(X)=0$, where $d\theta(X)=-J_3X^{\mathsf t}J_3$.
Conjugation by $1-X/2$ kills it. The remaining conjugations
by $1+Y$ with $d\theta(Y)=Y$ keep this mark. Their effect on
the difference at the second mark is $[Y,\bar\zeta T]$.
The prescribed two adjacent entries can be killed independently
using the reflected negative pair of entries in $Y$, because
$\bar\zeta$ is a unit. In each individual weight only the
corresponding correction is needed. The local conditions are
unchanged under these conjugations. The same normalization
argument for residual cocycles gives a normalized representative
of each class of $H^1_L(j)$. It follows that $h_j\le n_j$:
every such representative gives a tangent in the minimal
presentation, and the tangent space in weight $j$ surjects
linearly onto the indicated cohomology.

For the relation count, choose $s$ large enough that
$\mathfrak a\cap\mathfrak m_P^s\subset
\mathfrak m_P\mathfrak a$, by Artin--Rees. The kernel of
\[
 P/(\mathfrak m_P\mathfrak a+\mathfrak m_P^s)
          \longrightarrow P/(\mathfrak a+\mathfrak m_P^s)
\]
is then $\mathfrak a/\mathfrak m_P\mathfrak a$, killed by
$\mathfrak m_P$. These ideals are homogeneous, including after
truncation. Quotient out the other weight spaces of this
kernel. If its weight-$j$ dimension exceeded $b_j$, a nonzero
linear functional to $k$ would kill the obstruction in
$\operatorname{Ob}(j)\otimes_k D$.
Pushing the extension out along this functional would give a
normalized lift of the tautological point to the resulting
one-dimensional extension. Universality would give a map from
$R$ to this extension. Compare the induced map from $P$ with the
tautological quotient map: they differ on the variables only
by elements of the small kernel. Since all linear coefficients
of every relation are in $\pi O$ and its constant coefficient
is zero, its images under the two maps agree. But the
tautological map sends the relation space onto the nonzero
kernel, whereas the map through $R$ kills it. This contradiction
proves
\[
 \dim_k(\mathfrak a/\mathfrak m_P\mathfrak a)_j
       \le b_j\le n_j-(h_j-b_j).                               \tag{DC5}
\]
Nakayama gives actual homogeneous generators of $\mathfrak a$:
choose homogeneous lifts of a basis, using the weight projections
before passing to the finite quotient. Pad the lists with zero
equations to get respectively
\[
 (n_{-2},\ n_{-1}+1,\ n_0,\ n_1-1,\ n_2-1)
 \quad\text{in weights }(-2,-1,0,1,2).                         \tag{DC6}
\]
The two subtractions are legitimate by (DC4)--(DC5).
There are no equations of other weights needed.

\emph{The graded ring and the integral cotangent module.}
Let $A_{\rm gr}$ consist of the finite sums of homogeneous
formal series in $P$, with degree-zero ring $B$, and let $Q$
be the ideal in $A_{\rm gr}$ generated by the preceding equations.
The monoid of balanced exponent vectors has finitely many minimal
nonzero elements by Dickson's lemma. They generate it, by successive
subtraction. Thus $B$ is a quotient of a finite-variable formal
power series ring over $O$: choose a factorization of each balanced
monomial into the generators and substitute, which is legitimate
in the formal topology. In particular $B$ is Noetherian local
with residue field $k$. Each weight space is a finite sum of
minimal monomials times $B$, again by Dickson's lemma, so
$A_{\rm gr}$ is generated over $B$ by the variables.
Moreover
\[
 A_{\rm gr}/Q=\operatorname{Sring}(O,u),\qquad
 B/(Q\cap B)=R_0.                                               \tag{DC7}
\]
To see the kernel equality, express a homogeneous relation as a
finite sum of the homogeneous generators times arbitrary series
in $P$. Taking the component of the required weight expresses
it using homogeneous multipliers. For a finite sum of weights
repeat this operation finitely many times. The identification of
the image with the literal functional weight spaces is (DU1).
In particular no coefficient torsion is divided out in (DC7).
Theorem~\ref{thm:rh-dyadic-host-universal} gives finiteness of $R_0$
over $O$.

The augmentation $A_{\rm gr}\to O$ is evaluation at zero. Its
kernel $I$ is generated by the variables: partition each series
without constant term according to a variable dividing its monomials
and factor it out. Similarly $I^2$ consists exactly of the series
without constant or linear coefficients. Hence $I/I^2$ is free
over $O$ on the variables and $I^{(2)}=I^2$.
The ambient chart into $F[[z_1,\ldots,z_m]]$ is the literal inclusion.
The lists (DC6) and this chart are exactly those of
Theorem~\ref{thm:selmer-abstract-length}.

The original nonpositive-weight cotangent spanning holds.
Indeed a vector over $F$ annihilating the linear parts of the
equations in weight $j\le0$ can be multiplied by a nonzero
coefficient to have entries in $O$. It gives a normalized
deformation of the base point over $O[\delta]/(\delta^2)$.
Its cocycle has precisely weight $j$. The local equations give
the hypotheses of \texttt{VanE}: away from $\ell$ it kills
$N_v$, and at $\ell$ its weight-zero part vanishes on inertia
or its negative part is a coboundary on the multiplier-one
subgroup, as follows by linearizing the ordered flag.
After multiplying by a nonzero scalar, \texttt{VanE} makes it a
global coboundary. Over $F$, the infinitesimal conjugating matrix
$X$ for this coboundary satisfies $d\theta(X)=X$ at $c$.
By (DC3), the difference at the second mark is diagonal, so
$[X,\zeta T]$ is diagonal. The distinct eigenvalues imply
that $X$ is diagonal; $d\theta(X)=X$ then gives $X=\lambda T$.
Such a matrix centralizes the base polarized representation.
The cocycle and all entry displacements are zero.
Thus the original vector was zero, proving the asserted span.

\emph{The original characteristic-zero points.}
By (DC7) and the valuation argument (DU2), every graded point
over $\overline{\mathbf Q}_\ell$ scales by the torus into a
point of $R$ over a finite coefficient DVR. In particular its
matrix representation is continuous and satisfies the original
local equations. A point killing all nonzero weights is
diagonal. At $\ell$, the distinct residual eigenvalues of $x$
force its original ordered flag to be the standard flag. Its
diagonal inertia characters are therefore exactly the
prescribed ones. Away from $\ell$ the diagonal ratios to the
base point kill $N_v$ and on the remaining finite inertia
quotient have prime-to-$\ell$ order. Since they reduce to one,
they are trivial on this quotient. On $\ker\eta$ the three
ratios are now everywhere unramified characters, so \texttt{Cl1}
makes them trivial (apply it at each finite coefficient quotient).
The normalization at $c$ extends equality to $\Gamma$.
Because the variables are actual entries, the point is the
augmentation. In particular the fixed-point separation
hypothesis of the abstract length theorem holds.

Every two-sided point is either type I or type II. Work in its
integral torus conjugate. By (DC3) every invariant subspace
on $\ker\eta$ is a sum of coordinate lines, since the marked
difference is $\zeta T$ with three distinct eigenvalues.
Regard the nonzero matrix entries, for all group elements,
as edges directed from column to row. Reversing transpose
pairs the two upper adjacent edges and the two lower adjacent
edges. Both signs occur at a two-sided point, since the actual
entry variables generate the graded algebra over $R_0$.
If either adjacent pair occurs, an opposite-sign edge makes
the graph strongly connected: it is either a reversed adjacent
pair or the corner completing a cycle. No proper nonzero
coordinate subspace is invariant. Otherwise both adjacent
pairs vanish and both outer corners occur. The middle line
is separate and the outer block is irreducible, by its
two marked projectors and the nonzero entries in both
directions. These are exactly types I and II.

\emph{The rank-two arithmetic point and its rigidity.}
We retain the auxiliary field and quotient of
Theorem~\ref{thm:rh-dyadic-host-universal}, not the separate
good-data family. Write $A$ for the outer block on $G_K$
and $d$ for the middle character of a type-II point.
The Hensel basis at $x$ in (DH3) respects the blocks.
The ordered inertia characters of $A$ are
$\Psi,\Psi^{-1}$ and that of $d$ is one. Thus
$d,\det A$ are trivial on inertia at $\ell$.
Elsewhere they factor on inertia through the original
prime-to-$\ell$ quotient by $N_v$, and reduce to one,
so are trivial there as well. As $h_{\mathbf Q(i)}=1$,
both characters are one globally. The polarization gives
(RR3), hence the extension $W$ of (RR2) with determinant
$\eta$ and $W(c)=C=\operatorname{diag}(-1,1)$.
The strong irreducibility argument in
Lemma~\ref{lem:support-type-two-representation} applies:
its only additional inputs are the ordered local flag,
infinite-order $\Psi$ on inertia, and (DH3).
Indeed a reducible restriction to an open normal subgroup
is either scalar (excluded by infinite-order inertia)
or has two permuted lines. The local flag fixes both lines.
For a global element interchanging them, reducing its
trace and that after multiplication by $x$ gives
$\alpha+\alpha^{-1}=0$ and
$\alpha\beta+(\alpha\beta)^{-1}=0$, where
$\beta=\bar\Psi(x)$, contradicting $\beta^2\ne1$.
The connected monodromy group acts irreducibly and its
center acts by scalars, so its derived group acts
irreducibly; the cyclotomic kernel has Zariski closure
containing the derived group. This also works after
every finite extension.

Apply the actual normalized lift (FE1) to this inertial
point over its same coefficient DVR. Its outer block
has the two residual characters, the restricted marking
and polarization, and the exact ordered two-factor local
equations from its Hensel flag. This is the rank-two
normalized lift. The rank-two inputs to
Theorem~\ref{thm:support-finite-points} are supplied on
the very same auxiliary field: the host admission and
Proposition~\ref{prop:support-auxiliary-seed} give (AS1)
and (AF9), and Proposition~\ref{prop:selmer-arithmetic-seed}
constructs both ranks of the arithmetic seed.
The prescribed rank-two CM twist is supplied by
Theorem~\ref{thm:actual-cm-twist}.
The ordinary Hecke quotient, starting component and
rank-two propagation use the same finite-coefficient
hypotheses of Theorem~\ref{thm:support-auxiliary-frame}.
For the universal frame condition N4, the distinguished
element has scalar residual characteristic polynomial
and nontrivial residual cyclotomic value. The characteristic
polynomials of its matrix and of its cyclotomic multiple
are relatively prime modulo the coefficient maximal ideal.
Their resultant is a unit; over the fraction field
Cayley--Hamilton and coprimality kill every intertwiner
between those matrices. This checks N4 for every
original characteristic-prime point, as required,
without residual nonresonance at the coefficient prime.
In particular all these rank-two suppliers are on the
admitted host frame; the good-data end theorem is not
being applied to a resonant datum.

Strong irreducibility gives the rank-two field test.
The resulting actual Hecke point has the determinant
of $W|_{G_L}\kappa_2$ and the ordered local characters.
The Hodge labels differ by $2e>0$, so arithmetic-point
classicality makes the twist regular algebraic
conjugate self-dual cuspidal automorphic over $L$.
Now the untwisting and soluble descent argument of
Lemma~\ref{lem:support-type-two-newform} applies:
$L/\mathbf Q$ is abelian CM, the restriction to $G_L$
is irreducible, and $W$ has its odd determinant pairing.
BLGGT Lemmas 2.2.1 and 2.2.2 first untwist, then descend
to $\mathbf Q$. The Hodge labels and Frobenius
normalization (RR4) identify $W\varepsilon^{-e}$ with
the representation of a classical normalized cuspidal
newform of weight $2e+1$ and character $\eta$.
\cite{support-type-two-blggt,support-type-two-kato}

We check (SP1) over $\mathbf Q$, rather than just over
$K$. Away from $\ell$, inertia acts through the finite
abelian image of $(\eta,\psi)$ of modulus dividing $4N$.
At $\ell$, the ordered Hensel flag in the outer block is
invariant on the multiplier-one subgroup. If the
decomposition group contains a multiplier-minus-one
element, uniqueness of the marked ordered spectral flag
makes it preserve the flag in the polarized semidirect
sense, as in the local lifting argument above. In
dimension two
$\theta_2(h)=J_2h^{-\mathsf t}J_2=(\det h)^{-1}ChC$,
so in the outer Hensel basis $h$,
\[
 h^{-1}W(s)h=(\det h)
       \bigl(h^{-1}r(s)_{\rm outer}\theta_2(h)\bigr)C
\]
is upper triangular for every such $s$.
Inertia has multiplier one and its characters are
$\Psi,\Psi^{-1}$. Theorem~\ref{thm:rh-dyadic-host-crystalline}
makes the representation crystalline on restriction to
$K$, hence potentially crystalline over $\mathbf Q_\ell$.
This proves (SP1) with $M=4N$ and exponents $(e,-e)$.
Thus the type-II restriction belongs to the rank-two
class of Theorem~\ref{thm:support-spectrum}.

The adjoint calculation (RR7) uses just the polarized
outer block with determinant one, the trivial middle
line and $r(c)=1$, so over a finite coefficient field
$E$ it gives
\[
 \operatorname{ad}r=\operatorname{ad}^{0}W
          \oplus E(\eta)^{\oplus2}\oplus W\oplus W\eta .
\]
Thorne's Theorem 6.1 for the automorphic
$W|_{G_L}\kappa_2$, with its verified cyclotomic
irreducibility and polarization, kills the polarized
outer adjoint finite Selmer group. Restriction-corestriction
descends this to $\mathbf Q$; its scalar summand also
treats the second copy of $E(\eta)$.
For $W,W\eta$ use Kato's Theorem 14.2(1) over the
abelian field $K$ for the newform at the critical
integer $e$: $1\le e\le2e$ and $e\ne(2e+1)/2$.
Lemma~\ref{lem:support-type-two-kato} proves the
passage from his finite discrete group to rational
vanishing, with the same infinite-order inertia element.
Restriction-corestriction gives the two vanishings over
$\mathbf Q$. Therefore $H_f^1(\mathbf Q,\operatorname{ad}r)=0$.
All original inertial tangents are in this finite group
by Theorem~\ref{thm:rh-dyadic-host-crystalline} and the
away-from-$\ell$ averaging argument of
Theorem~\ref{thm:rh-dyadic-host-traces}, neither of which
requires global irreducibility.

The type-I points were treated in
Theorems~\ref{thm:rh-dyadic-host-universal}
and~\ref{thm:rh-dyadic-host-traces}. Both types now
belong to the classes of Theorem~\ref{thm:support-spectrum}.
Using exactly the Frobenius chosen in (DH6), (SP5)--(SP6)
give (DH6) for every mixed family of the two types,
and nonvanishing at each two-sided point.
No weight-minus-two vanishing is needed.

\emph{Tangents and the trace element.}
For a two-sided point $\xi$ and any derivation at it, first
apply the same torus scaling as for the point itself.
The values of the derivation on the finite $O$-module $R_0$
and on the finitely many variables lie in one finite extension
of $F$. Multiplying the derivation by a nonzero scalar makes
these values integral. The map $\xi+\delta D$ on $R_0$ is
continuous, since it is $O$-linear on a finite module; its
residue is the same as that of $\xi$. Each variable value
belongs to the maximal ideal of $O'[\delta]/(\delta^2)$.
The decomposition into minimal monomials times degree-zero
series in (DU2) now shows that this dual-number map extends by
convergent evaluation to $R$. Indeed it agrees with the
given map on every homogeneous part, so every homogeneous
part of every relation vanishes; the sum converges to zero.
The resulting continuous deformation satisfies the original
local equations. The rigidity arguments for the two types make its tangent
a coboundary over $\overline{\mathbf Q}_\ell$.
The calculation with the two marks above gives
$X=\lambda T$ again, now at the given point. It follows that
on every matrix entry, and then on every homogeneous series
by the evaluation argument, its derivation is
\[
                  D(f)=\lambda j\,\xi(f)\quad(f\text{ of weight }j).
\]
Undoing the scalings preserves this assertion. This is exactly
the required one-dimensional tangent bound, including
nilpotent tangents.

Take as $a$ in the abstract length theorem a degree-zero
representative of
$\operatorname{tr}u(\phi)-\operatorname{tr}r_0(\phi)$
in $B$, which exists by (DC7) and has augmentation zero.
Its value is nonzero at every two-sided point, by the
automorphy of those points and (DH6). The hypotheses of the
abstract length theorem are now all verified, with $r=n_{-1}$
in its notation. If $n_{-1}=0$, the length comparison below
is zero and needs no application of that theorem.

\emph{Comparison at every modulus.}
Write $\mathcal C^-$ for the weight-$-1$ quotient of $I/I^2$
by the linear parts of all weight-$-1$ relations. For
$0\ne a_0\in O$ take a class in the original
$\operatorname{Sha}_O(a_0)$, with coefficients
$(O/a_0)(\Psi^{-1})$. The set of classes represented by
cocycles with the required local conditions is itself an
$O$-submodule, so this also describes the span in the original
definition. There is a unique representative $z$ satisfying
$z(g_1)=0$: add a scalar coboundary and use that
$\Psi(g_1)^{-1}-1$ is a unit. The local conditions persist,
since $\Psi(N_v)=1$ and the condition at $\ell$ is being a
coboundary. As $2$ is invertible, $z(c)=0$; also
$z(g_1^*)=0$ by the cocycle identity.

The adjoint action on $L=E_{10}-E_{21}$ is multiplication by
$\Psi^{-1}$, including on the other coset because
$d\theta(L)=L$. The matrices
\[
                     (1+\delta z(g)L)r_0(g)
\]
over the square-zero ring $O\oplus (O/a_0)\delta$ therefore
give the actual normalized deformation. At $\ell$ its
coboundary is effected by conjugation of the base point by a
matrix congruent to one; away from $\ell$ it kills $N_v$.
Continuity is inherited from the original cocycle.
Evaluating the actual universal variables gives an $O$-linear
map
\[
            \operatorname{Sha}_O(a_0)
                  \longrightarrow
                 \operatorname{Hom}_O(\mathcal C^-,O/a_0).      \tag{DC9}
\]
The evaluation is precisely a first-order displacement in
weight $-1$, and all relation linear parts vanish at it.
If the displacement vector is zero, the universal matrices
are the base matrices, so $z=0$. The normalized representative
depends linearly on the class, proving injectivity.
This argument is valid at every modulus, with no passage to
a saturated quotient. If $n_{-1}=0$, take $V=0$ and an empty
family in (DC1). Otherwise take
$V=\operatorname{length}_O\mathcal C^-$, finite by the
nonpositive span. By the elementary divisors of this finite
module, the length of the right side of (DC9) is at most $V$
for every $a_0$.

Apply Theorem~\ref{thm:selmer-abstract-length} to the preceding
graded ring, with the finite set of its two-sided torus-orbit
representatives. Their representations are pairwise
nonisomorphic on $\ker\eta$. Indeed an isomorphism intertwines
the two matrices $\zeta T,\zeta'T$ of (DC3); their spectra
force $\zeta'=\pm\zeta$, and their common nonzero residual
value excludes the negative sign. The intertwiner is thus
diagonal. If it is $M$, in type I the reversing-transpose identity
for $g^*$ shows that $M J_3M^{\mathsf t}J_3$ commutes with
the irreducible matrix representation. By Schur's lemma it
is scalar. For diagonal $M$ this says
$M_{00}M_{22}=M_{11}^2$, so, up to a scalar, $M$ is
$\operatorname{diag}(t,1,t^{-1})$.
In type II choose $t$ with $t^2=M_{00}/M_{22}$ in the
algebraic closure; the torus has exactly the same action
on the outer block and fixes the trivial middle line.
Points of different types cannot be isomorphic.
Both matrices at $c$ are one. The two points are therefore torus conjugate on the
whole group and on all graded series, by their integral
evaluation on the actual entry variables. Distinct orbits
cannot give such an isomorphism.

The abstract length bound now gives (DC1), since
$v_\pi=e_Ov_\ell$. The strict lower comparison of
Theorem~\ref{thm:rh-dyadic-host-frame} bounds the original
strict length by $V$ (equivalently, apply (DC9) to each finite
torsion subgroup). Finally divide by $e_O$ and apply (DH6).
Absorbing the factor two into an absolute exponent gives
(DC2). This proves the integral estimate for the exact
inertial ring, with its original equations and all its
coefficient torsion.
\end{proof}

